\beginsong{În trecerea grăbită prin lumea către veci}[by={Traian Dorz și Nicolae Moldoveanu}]
\beginverse
În trecerea grăbită prin lumea către veci,
Fă-ți timp măcar o clipă să vezi pe unde mergi!
Fă-ți timp să vezi durerea și lacrima arzând,
Fă-ți timp să poți cu milă să le alini trecând;
Fă-ți timp pentru-adevăruri și adânciri în vis;
Fă-ți timp pentru cântare cu sufletul deschis;
Fă-ți timp să vezi pădurea, s-asculți lângă izvor,
Fă-ți timp s-asculți ce spune o floare, un cocor;
Fă-ți timp s-aștepți din urmă când mergi cu slăbănogi,
Fă-ți timp pe-un munte seara stând singur să te rogi;
Fă-ți timp să stai cu mama și tatăl tău bătrâni,
Fă-ți timp de-o vorbă bună și-o coajă pentru câini.
/: Fă-ți timp, fă-ți timp, o, nu uita, fă-ți timp! :/
\endverse

\beginverse
În trecerea grăbită prin lume către veci,
Fă-ți timp măcar o clipă să vezi pe unde treci!
Fă-ți timp să stai de-aproape cu-ai tăi iubiți voios,
Fă-ți timp să fii și-al casei și-n slujba lui Hristos;
Fă-ți timp să guști frumsețea din tot ce e curat,
Fă-ți timp, că ești de taine și lumi înconjurat;
Fă-ți timp cu orice taină sau adevăr să stai,
Fă-ți timp, căci toate-acestea au inimă și grai;
Fă-ți timp s-asculți la toate, din toate să înveți,
Fă-ți timp să dai vieții adevăratul preț!
Fă-ți timp acum, că-n urmă zadarnic ai să plângi,
Comoara risipită a vieții n-o mai strângi!
/: Fă-ți timp, fă-ți timp, o, nu uita, fă-ți timp! :/
\endverse
\endsong